\documentclass[11pt,a4paper]{article}
\usepackage[margin=1in]{geometry}
\usepackage[T1]{fontenc}
\usepackage[utf8]{inputenc}
\usepackage{lmodern}
\usepackage{microtype}
\usepackage{hyperref}
\usepackage{enumitem}
\usepackage{xcolor}
\usepackage{longtable}
\usepackage{booktabs}
\usepackage{array}
\usepackage{graphicx}

\hypersetup{
  colorlinks=true,
  linkcolor=blue!60!black,
  urlcolor=blue!60!black
}

\newcommand{\gdshot}[3]{
  \begin{figure}[h]
    \centering
    \IfFileExists{#1}{
      \includegraphics[width=#2]{#1}
    }{
      \fbox{
        \parbox{0.85\linewidth}{
          \centering
          Screenshot placeholder\\[0.2em]
          \texttt{#1}
        }
      }
    }
    \caption{#3}
  \end{figure}
}

\newcommand{\gdicon}[1]{\includegraphics[height=1.8ex]{#1}}

\title{\textbf{GeoDraw User Guide}\\\large End-User Manual}
\author{GeoDraw}
\date{\today}

\begin{document}
\begin{titlepage}
  \centering
  {\Huge\bfseries GeoDraw User Guide\par}
  \vspace{0.6em}
  {\Large End-User Manual\par}
  \vspace{2em}
  {\large Version: Stable Desktop Build\par}
  \vfill
  {\large \today\par}
\end{titlepage}
\tableofcontents
\bigskip

\section{Purpose}
GeoDraw is an interactive geometry app for constructing points, lines, circles, angles, sectors, and polygons, then exporting the figure to TikZ (\texttt{tkz-euclide} style workflow).

\section{Interface Overview}
The window is divided into:
\begin{itemize}[leftmargin=1.5em]
  \item \textbf{Left Toolbar}: tool groups (Move, Points, Lines, Angle, Circles, Styles).
  \item \textbf{Center Canvas}: drawing and interaction area.
  \item \textbf{Top Actions}: file controls (Open, Save, Save As), Undo/Redo, Fit View.
  \item \textbf{Right Sidebar}: \textbf{Objects} tab and \textbf{Export} tab.
  \item \textbf{Bottom Command Bar}: expression evaluation and command-driven construction.
\end{itemize}

% \begin{figure}[h]
% \centering
% \includegraphics[width=\linewidth]{screenshots/layout-overview.png}
% \caption{Main layout}
% \end{figure}

\section{File Operations}
\subsection{Open / Save / Save As}
\begin{itemize}[leftmargin=1.5em]
  \item \textbf{Open}: load \texttt{.geodraw} or \texttt{.json} snapshot.
  \item \textbf{Save}: overwrite current file if already opened/saved once.
  \item \textbf{Save As}: choose new name/location.
\end{itemize}

\subsection{Supported file types}
\begin{itemize}[leftmargin=1.5em]
  \item \texttt{.geodraw}
  \item \texttt{.json} (same snapshot structure)
\end{itemize}

\section{Navigation and Global Controls}
\subsection{Mouse}
\begin{itemize}[leftmargin=1.5em]
  \item Drag points to move them.
  \item Pan the canvas in Move mode.
  \item Use mouse wheel/trackpad to zoom.
\end{itemize}

\subsection{Keyboard shortcuts}
\begin{itemize}[leftmargin=1.5em]
  \item \texttt{Ctrl/Cmd + Z}: Undo
  \item \texttt{Ctrl/Cmd + Shift + Z} or \texttt{Ctrl/Cmd + Y}: Redo
  \item \texttt{Shift + F}: Fit View
  \item \texttt{Delete}/\texttt{Backspace}: delete selected object
  \item \texttt{Esc}: cancel Copy Style mode / clear pending selections
\end{itemize}

\subsection{Fit View}
Use \textbf{Fit View} \gdicon{screenshots/icons/fit-view.png} from the top action row (next to Undo/Redo), or press \texttt{Shift + F}.
\begin{itemize}[leftmargin=1.5em]
  \item Re-centers the camera to include current scene geometry.
  \item Applies padding so objects are not touching viewport edges.
  \item If no geometry is available, camera stays unchanged.
\end{itemize}
\gdshot{screenshots/fit-view.png}{0.9\linewidth}{Fit View action and resulting re-centered camera.}

\section{Tool Palette}
Long-press or right-click a group button to open its flyout tools.

\subsection{Move Group}
\begin{itemize}[leftmargin=1.5em]
  \item \gdicon{screenshots/tools/move.png} \textbf{Move / Select} (\texttt{V}): select and manipulate objects.
  \item \gdicon{screenshots/tools/export_clip_rect.png} \textbf{Export Clip Rectangle}: create rectangular clip for export.
  \item \gdicon{screenshots/tools/export_clip.png} \textbf{Export Clip Polygon}: create polygon clip for export (close by clicking near the first vertex).
\end{itemize}

\subsection{Points Group}
\begin{itemize}[leftmargin=1.5em]
  \item \gdicon{screenshots/tools/point.png} \textbf{Point} (\texttt{P}): create/select points.
  \item \gdicon{screenshots/tools/midpoint.png} \textbf{Midpoint} (\texttt{M}): midpoint from two points or from a segment.
\end{itemize}

\subsection{Transform Group}
\begin{itemize}[leftmargin=1.5em]
  \item \gdicon{screenshots/tools/translate.png} \textbf{Translate}: create translated copy from vector \texttt{AB}.
  \item \gdicon{screenshots/tools/rotate.png} \textbf{Rotate}: rotate object around selected center.
  \item \gdicon{screenshots/tools/reflect.png} \textbf{Reflect}: reflect object across line/segment/point.
  \item \gdicon{screenshots/tools/dilate.png} \textbf{Dilate}: scale object from selected center.
\end{itemize}

\subsection{Lines Group}
\begin{itemize}[leftmargin=1.5em]
  \item \gdicon{screenshots/tools/segment.png} \textbf{Segment} (\texttt{S})
  \item \gdicon{screenshots/tools/line2p.png} \textbf{Line Through 2 Points} (\texttt{L})
  \item \gdicon{screenshots/tools/perp_line.png} \textbf{Perpendicular Line}
  \item \gdicon{screenshots/tools/parallel_line.png} \textbf{Parallel Line}
  \item \gdicon{screenshots/tools/tangent_line.png} \textbf{Tangent Line}
  \item \gdicon{screenshots/tools/angle_bisector.png} \textbf{Internal Angle Bisector}
\end{itemize}

\subsection{Angle Group}
\begin{itemize}[leftmargin=1.5em]
  \item \gdicon{screenshots/tools/angle.png} \textbf{Angle}: select 3 points (A, vertex B, C).
  \item \gdicon{screenshots/tools/angle_fixed.png} \textbf{Angle with Fixed Value}: define expression and direction (CW/CCW), then construct.
\end{itemize}

\subsection{Circles Group}
\begin{itemize}[leftmargin=1.5em]
  \item \gdicon{screenshots/tools/circle_cp.png} \textbf{Circle Center + Point} (\texttt{O})
  \item \gdicon{screenshots/tools/circle_3p.png} \textbf{Circle through 3 Points}
  \item \gdicon{screenshots/tools/circle_fixed.png} \textbf{Circle with Fixed Radius}
  \item \gdicon{screenshots/tools/sector.png} \textbf{Circular Sector}
  \item \gdicon{screenshots/tools/polygon.png} \textbf{Polygon}
  \item \gdicon{screenshots/tools/regular_polygon.png} \textbf{Regular Polygon}
\end{itemize}

\subsection{Styles Group}
\begin{itemize}[leftmargin=1.5em]
  \item \gdicon{screenshots/tools/copyStyle.png} \textbf{Copy Style} (\texttt{C}): pick a source style, then apply to targets.
  \item \gdicon{screenshots/tools/label.png} \textbf{Label Tool}: place text labels on the canvas.
\end{itemize}

\section{Objects Tab (Right Sidebar)}
\subsection{Top toggles}
\begin{itemize}[leftmargin=1.5em]
  \item \textbf{Grid}: show/hide grid.
  \item \textbf{Axes}: show/hide axes.
  \item \textbf{Snap}: snap free-point creation to grid intersections.
  \item \textbf{Dependency Glow}: highlight dependency categories visually.
\end{itemize}

\subsection{Object Browser}
\begin{itemize}[leftmargin=1.5em]
  \item Tabs: All, Points, Lines, Circles, Angles, Numbers.
  \item Click row to select object.
  \item Use the row checkbox (right side) to toggle object visibility.
\end{itemize}

\section{Properties Panel}
When an object is selected, the panel shows editable properties.

\subsection{Point properties}
\begin{itemize}[leftmargin=1.5em]
  \item Rename point.
  \item Position readout.
  \item Visual style: shape, size, stroke/fill, label style.
\end{itemize}

\subsection{Line/Segment properties}
\begin{itemize}[leftmargin=1.5em]
  \item Stroke color, width, dash, opacity.
  \item Segment marks (\texttt{|}, \texttt{||}, \texttt{|||}, etc.) and arrow marks (end/mid).
\end{itemize}

\subsection{Arrow controls}
Arrow styling \gdicon{screenshots/icons/arrow-controls.png} is available in Segment, Circle, and Angle (arc) style panels.
\begin{itemize}[leftmargin=1.5em]
  \item Multi-arrow list editor: add, duplicate, remove, and select arrow entries.
  \item Direction options: forward/backward and paired directions.
  \item Tip style: Stealth, Latex, Triangle.
  \item Placement and distribution controls (single/multi, plus segment end/mid modes where applicable).
  \item Appearance controls: color, width, size, and arrow length.
\end{itemize}
\gdshot{screenshots/arrow-controls.png}{0.78\linewidth}{Arrow list editor in object style properties.}

\subsection{Circle/Polygon/Sector properties}
\begin{itemize}[leftmargin=1.5em]
  \item Stroke and fill settings.
  \item Fill opacity.
  \item Fill pattern and pattern color.
\end{itemize}

\subsection{Angle properties}
\begin{itemize}[leftmargin=1.5em]
  \item Arc radius, stroke, fill.
  \item Label/value visibility and placement.
  \item Mark styles:
  \begin{itemize}[leftmargin=1.5em]
    \item Non-right: vanilla arc, bar marks (\texttt{|}, \texttt{||}, \texttt{|||}), double/triple arc, none.
    \item Right-angle status aware options for exact/proven right angles.
  \end{itemize}
\end{itemize}

\subsection{Marking controls}
Marking tools \gdicon{screenshots/icons/marking-controls.png} support both segment congruence marks and angle arc marks.
\begin{itemize}[leftmargin=1.5em]
  \item \textbf{Segment marks}: symbol, distribution, position, size, width, and color.
  \item \textbf{Angle marks}: single/double/triple arc count, bar symbol, position, size, and color.
  \item Both use list-based editors, so multiple mark entries can be layered.
\end{itemize}
\gdshot{screenshots/mark-controls.png}{0.78\linewidth}{Segment and angle marking controls in the properties panel.}

\subsection{Tool info panels}
Special tools (Fixed Angle, Fixed Radius Circle, Export Clip, Copy Style) show context-specific instructions and inputs.

\section{Command Bar}
Command Bar supports expression evaluation, object creation commands, and assignments.

\subsection{Input behavior}
\begin{itemize}[leftmargin=1.5em]
  \item \texttt{Enter}: run command
  \item \texttt{Esc}: clear input
  \item \texttt{Up/Down}: command history
\end{itemize}

\subsection{Expression examples}
\begin{itemize}[leftmargin=1.5em]
  \item \texttt{5*5}
  \item \texttt{sin(pi/6)}, \texttt{atan2d(4,3)}
  \item \texttt{Distance(A,B)}
\end{itemize}

\subsection{Supported construction commands}
\begin{itemize}[leftmargin=1.5em]
  \item \textbf{Point/transform}: \texttt{Point(x,y)}, \texttt{Midpoint(A,B)}, \texttt{Midpoint(s)}, \texttt{Incenter(A,B,C)}, \texttt{Orthocenter(A,B,C)} (\texttt{Ortho}), \texttt{Centroid(A,B,C)}, \texttt{Translate(P,A,B)}, \texttt{Rotate(P,O,expr[,CW|CCW])}, \texttt{Dilate(P,O,k)}, \texttt{Reflect(P,l|s|O)}.
  \item \textbf{Line/segment}: \texttt{Line(A,B)}, \texttt{Line(x1,y1,x2,y2)}, \texttt{Perpendicular(P,l)}, \texttt{Parallel(P,l)}, \texttt{Tangent(P,c)}, \texttt{AngleBisector(A,B,C)}, \texttt{Segment(A,B)}.
  \item \textbf{Angle/polygon}: \texttt{Angle(A,B,C)}, \texttt{AngleFixed(V,A,expr[,CW|CCW])}, \texttt{Sector(O,A,B)}, \texttt{Polygon(A,B,C,...)}, \texttt{RegularPolygon(A,B,n[,CW|CCW])}.
  \item \textbf{Circle}: \texttt{Circle(x,y,r)}, \texttt{Circle(O,A)}, \texttt{Circle(O,r)}, \texttt{Circle3P(A,B,C)} (alias: \texttt{CircleThreePoint(A,B,C)}).
  \item \textbf{Numeric query}: \texttt{Distance(...)} returns a scalar value (no object created).
\end{itemize}

\subsection{Assignments}
\begin{itemize}[leftmargin=1.5em]
  \item Scalar: \texttt{n\_1 = 2.023}
  \item Named object: \texttt{P = Point(1,2)}, \texttt{l = Line(A,B)}, \texttt{c = Circle(O,r\_1)}
  \item Point-expression assignment (vector style): \texttt{X = A + B}
\end{itemize}

\subsection{Important assignment rule}
Assignment is selective and fail-closed:
\begin{itemize}[leftmargin=1.5em]
  \item Existing constant numbers are updated in-place by scalar assignment.
  \item Existing free points are updated in-place by point assignment.
  \item Existing aliases can be redefined only with compatible constructor types.
  \item Non-free points, non-constant numbers, and incompatible alias redefines are rejected.
\end{itemize}

\section{Export Tab}
\subsection{TikZ export options}
\begin{itemize}[leftmargin=1.5em]
  \item \textbf{Use Current View}: export based on current camera viewport.
  \item \textbf{Use Export Clip selection}: apply clip polygon from Export Clip tool.
  \item \textbf{Match canvas size conversion}
  \item \textbf{Label glow}
  \item Scale controls:
  \begin{itemize}[leftmargin=1.5em]
    \item Global Scale
    \item Point Scale
    \item Line Scale
    \item Label Scale
  \end{itemize}
\end{itemize}

\subsection{Preview window (PDF + editable TikZ)}
Press \textbf{Preview} \gdicon{screenshots/icons/pdf-preview.png} in Export to open the dedicated TikZ preview window.
\begin{itemize}[leftmargin=1.5em]
  \item Left pane: compiled PDF preview from the current TikZ source.
  \item Right pane: editable \texttt{tikzpicture} code.
  \item \textbf{Optional Preamble} section for custom preamble lines (for example \texttt{\textbackslash pagecolor\{...\}}).
  \item Find/replace tools with case-match option.
  \item \textbf{Update PDF} recompiles, and \textbf{Copy Edited TikZ} copies the edited code.
  \item Compiler log panel is available for diagnostics when compilation fails.
\end{itemize}
\gdshot{screenshots/tikz-preview-window.png}{0.95\linewidth}{TikZ Preview window with PDF pane, code editor, and compiler log.}

\subsection{JSON export}
\begin{itemize}[leftmargin=1.5em]
  \item Export model snapshot as JSON.
  \item Optional debug world coordinates.
\end{itemize}

\section{Typical Workflows}
\subsection{Workflow A: classic construction}
\begin{enumerate}[leftmargin=1.5em]
  \item Create base points with Point tool.
  \item Build segments/lines/circles.
  \item Add angles/sectors/polygons.
  \item Tune style from Properties panel.
  \item Export TikZ.
\end{enumerate}

\subsection{Workflow B: command-driven}
\begin{enumerate}[leftmargin=1.5em]
  \item Create anchor points from Command Bar.
  \item Create lines/circles/segments via named references.
  \item Assign constants and reuse them in commands.
  \item Export.
\end{enumerate}

\section{Troubleshooting}
\subsection{I cannot select a tool variant}
Long-press or right-click the group button to open its flyout.

\subsection{Clip export does not apply}
Ensure you closed the clip polygon (click near first point), then enable \textbf{Use Export Clip selection} in Export tab.

\subsection{Command reports unknown/ambiguous point}
\begin{itemize}[leftmargin=1.5em]
  \item Verify point labels exist.
  \item Avoid duplicate labels.
\end{itemize}

\subsection{Object is hidden unexpectedly}
Check per-object visibility checkbox in Object Browser.

\section{Current Limits (User-visible)}
\begin{itemize}[leftmargin=1.5em]
  \item Tangent assignment is intentionally blocked because \texttt{Tangent(P,c)} may create multiple lines.
  \item Alias-based commands (\texttt{Midpoint(s)}, \texttt{Perpendicular}, \texttt{Parallel}, \texttt{Tangent}) require prior named aliases.
  \item Some advanced symbolic redefine workflows are intentionally fail-closed.
\end{itemize}

\appendix
\section{Screenshot Appendix Template}
Use this template to insert screenshots consistently.

\subsection*{A.1 Main Layout}
% \begin{figure}[h]
% \centering
% \includegraphics[width=\linewidth]{screenshots/main-layout.png}
% \caption{Main layout with tool palette, canvas, sidebar, and command bar.}
% \end{figure}

\subsection*{A.2 Tool Flyout}
% \begin{figure}[h]
% \centering
% \includegraphics[width=.7\linewidth]{screenshots/tool-flyout.png}
% \caption{Long-press/right-click flyout for tool variants.}
% \end{figure}

\subsection*{A.3 Object Browser}
% \begin{figure}[h]
% \centering
% \includegraphics[width=.7\linewidth]{screenshots/object-browser.png}
% \caption{Object tabs and per-object visibility checkbox.}
% \end{figure}

\subsection*{A.4 Properties Panel}
% \begin{figure}[h]
% \centering
% \includegraphics[width=.7\linewidth]{screenshots/properties-panel.png}
% \caption{Example properties for a selected object.}
% \end{figure}

\subsection*{A.5 Command Bar}
% \begin{figure}[h]
% \centering
% \includegraphics[width=\linewidth]{screenshots/command-bar.png}
% \caption{Command entry and status feedback.}
% \end{figure}

\subsection*{A.6 Export Panel + Clip}
% \begin{figure}[h]
% \centering
% \includegraphics[width=\linewidth]{screenshots/export-panel.png}
% \caption{Export options and clip selection workflow.}
% \end{figure}

\subsection*{A.7 TikZ Output}
% \begin{figure}[h]
% \centering
% \includegraphics[width=\linewidth]{screenshots/tikz-output.png}
% \caption{Compiled TikZ result preview.}
% \end{figure}

\end{document}
